\documentclass[11pt,a4paper]{article}
\usepackage[utf8]{inputenc}
\usepackage{amsmath,amssymb,amsthm}
\usepackage{mathtools}
\usepackage{geometry}
\usepackage{CJKutf8}
\usepackage{fancyhdr}
\usepackage{tcolorbox}
\usepackage{xcolor}
\usepackage{enumitem}

\geometry{margin=2.5cm, top=3.5cm, bottom=2.5cm, headheight=28pt}

% ========== 页眉设置 ==========
\pagestyle{fancy}
\fancyhf{}
\fancyhead[L]{\textsc{Machine Learning Theory}\\Lecture 3}
\fancyhead[R]{\textsc{Machine Learning Theory Lecture 3}\\Generalized Interpolating Discrete Diffusion}
\fancyfoot[C]{\thepage}
\renewcommand{\headrulewidth}{0pt}

% ========== tcolorbox 环境定义 ==========
\tcbuselibrary{theorems, skins, breakable}

% 定义框（蓝色边框）
\newtcolorbox{defbox}[1][]{
  enhanced,
  breakable,
  colback=white,
  colframe=blue!70!black,
  boxrule=0.8pt,
  left=8pt, right=8pt, top=6pt, bottom=6pt,
  before skip=10pt, after skip=10pt,
  #1
}

% 定理框（蓝色边框）
\newtcolorbox{thmbox}[1][]{
  enhanced,
  breakable,
  colback=white,
  colframe=blue!70!black,
  boxrule=0.8pt,
  left=8pt, right=8pt, top=6pt, bottom=6pt,
  before skip=10pt, after skip=10pt,
  #1
}

% 引理框
\newtcolorbox{lembox}[1][]{
  enhanced,
  breakable,
  colback=white,
  colframe=blue!70!black,
  boxrule=0.8pt,
  left=8pt, right=8pt, top=6pt, bottom=6pt,
  before skip=10pt, after skip=10pt,
  #1
}

% 命题框
\newtcolorbox{propbox}[1][]{
  enhanced,
  breakable,
  colback=white,
  colframe=blue!70!black,
  boxrule=0.8pt,
  left=8pt, right=8pt, top=6pt, bottom=6pt,
  before skip=10pt, after skip=10pt,
  #1
}

% 备注框（灰色背景）
\newtcolorbox{remarkbox}[1][]{
  enhanced,
  breakable,
  colback=gray!10,
  colframe=gray!10,
  boxrule=0pt,
  left=8pt, right=8pt, top=6pt, bottom=6pt,
  before skip=10pt, after skip=10pt,
  #1
}

% ========== 计数器 ==========
\newcounter{proposition}[section]
\newcounter{lemma}[section]
\newcounter{theorem}[section]
\newcounter{definition}[section]
\newcounter{remark}[section]

\renewcommand{\theproposition}{\thesection.\arabic{proposition}}
\renewcommand{\thelemma}{\thesection.\arabic{lemma}}
\renewcommand{\thetheorem}{\thesection.\arabic{theorem}}
\renewcommand{\thedefinition}{\thesection.\arabic{definition}}
\renewcommand{\theremark}{\thesection.\arabic{remark}}

% ========== 子章节格式 ==========
\renewcommand{\thesubsection}{\fbox{\thesection.\arabic{subsection}}}

% ========== 数学算子 ==========
\DeclareMathOperator{\Cat}{Cat}
\newcommand{\Dkl}{D_{\mathrm{KL}}}
\newcommand{\Dis}{D_{\mathrm{IS}}}

\begin{document}
\begin{CJK}{UTF8}{gbsn}

% ========== 标题 ==========
\begin{center}
{\LARGE\bfseries 1\quad GIDD: Generalized Interpolating Discrete Diffusion}
\end{center}

\setcounter{section}{1}

% ========== 1.1 基本定义 ==========
\subsection{基本定义}

\refstepcounter{definition}
\begin{defbox}
\textbf{定义 \thedefinition（Mixing Rate）.}
\textbf{累积混合率}（Cumulative Mixing Rate）定义为可微的单调递减函数
\[
\alpha: [0,1] \to [0,1], \quad t \mapsto \alpha_t
\]
满足边界条件 $\alpha_0 = 1$（无混合）且 $\alpha_1 = 0$（完全混合）。

定义 $\beta_t := 1 - \alpha_t$，则信噪比（SNR）为
\[
\mathrm{SNR}(t) := \frac{\alpha_t}{\beta_t}
\]
\end{defbox}

\refstepcounter{definition}
\begin{defbox}
\textbf{定义 \thedefinition（Mixing Distribution）.}
设 $|V|$ 为词表大小。\textbf{混合分布}（Mixing Distribution）定义为可微函数
\[
\bar{\pi}: [0,1] \to \Delta^{|V|-1}, \quad t \mapsto \bar{\pi}_t
\]
其中 $\bar{\pi}_t := \bar{\pi}(t)$ 表示时刻 $t$ 的混合分布，$\Delta^{d-1}$ 为 $d$ 维单纯形：
\[
\Delta^{d-1} = \left\{ \mathbf{x} \in \mathbb{R}^d \;\middle|\; x_i \geq 0,\; \sum_{i=1}^{d} x_i = 1 \right\}
\]

分布 $\bar{\pi}_t$ 决定了在任意时刻 $t$ 添加到数据的噪声类型。$\bar{\pi}_1$ 代表扩散过程的先验分布。
\end{defbox}

\refstepcounter{definition}
\begin{defbox}
\textbf{定义 \thedefinition（Marginal Forward Transition）.}
设 $x \in V$ 为数据点，$\mathbf{x} \in \{0,1\}^{|V|}$ 为其 one-hot 编码。\textbf{边际前向转移}定义为
\[
q_t(z_t|x) := \Cat\left(z_t;\, \alpha_t \mathbf{x} + \beta_t \bar{\pi}_t\right)
\]
其中 $z_t \in V$ 为时刻 $t$ 的隐变量。

当 $\bar{\pi}_t = \mathbf{m}$（mask token 的 one-hot 向量）时，GIDD 退化为 MDM（Masked Diffusion Model）。
\end{defbox}

% ========== 1.2 构造连续时间 Markov 矩阵 ==========
\subsection{构造连续时间 Markov 矩阵}

\refstepcounter{proposition}
\begin{propbox}
\textbf{命题 \theproposition.}
存在连续时间 Markov 链，从状态 $z_s$ 到 $z_t$ ($s \leq t$)  的转移概率为
\[
q_{t|s}(z_t|z_s) = \Cat(z_t; Q_{t|s} z_s), \quad Q_{t|s} = \alpha_{t|s} I + \beta_{t|s} \bar{\pi}_{t|s} \mathbf{1}^T
\]
其中 $\alpha_{t|s} = \dfrac{\alpha_t}{\alpha_s}$，$\beta_{t|s} \bar{\pi}_{t|s} = \beta_t \bar{\pi}_t - \dfrac{\alpha_t}{\alpha_s} \beta_s \bar{\pi}_s$，$\mathbf{1}$ 为 $|V|$ 维分量全为 1 的向量。
\end{propbox}

% ========== 1.3 命题1.1的证明 ==========
\subsection{命题1.1的证明}

对任意给定的 $\Delta > 0$，将 $[0,1]$ 等距划分为网格 $\{0, \Delta, 2\Delta, \ldots, 1\}$。对于 $i \in \mathbb{Z}_{\geq 0}$，定义：
\[
\dot{\alpha}_{\Delta i} := \frac{\alpha_{\Delta(i+1)}}{\alpha_{\Delta i}}, \quad
\dot{\beta}_{\Delta i} \dot{\bar{\pi}}_{\Delta i} := \beta_{\Delta(i+1)} \bar{\pi}_{\Delta(i+1)} - \frac{\alpha_{\Delta(i+1)}}{\alpha_{\Delta i}} \beta_{\Delta i} \bar{\pi}_{\Delta i}
\]

定义单步转移概率：
\[
q_t(z_{t+\Delta}|z_t) = \Cat(z_{t+\Delta}; \dot{Q}_t z_t), \quad \dot{Q}_t = \dot{\alpha}_t I + \dot{\beta}_t \dot{\bar{\pi}}_t \mathbf{1}^T
\]
这是一个离散时间 Markov 链。对于 $s, t \in \{0, \Delta, 2\Delta, \ldots, 1\}$ 且 $s \leq t$，从时刻 $s$ 到时刻 $t$ 的累积转移矩阵为
\[
Q_{t|s} = \prod_{i=s/\Delta}^{\frac{t}{\Delta}-1} \dot{Q}_{\Delta i}
\]

\textbf{用数学归纳法证明}

\[
Q_{t|s} = \alpha_{t|s} I + \beta_{t|s} \bar{\pi}_{t|s} \mathbf{1}^T
\]

\textbf{证明.}

\textbf{① 当 $t=s$ 时，}
\[
Q_{t|s} = Q_{s|s} = \alpha_{s|s} I + \beta_{s|s} \bar{\pi}_{s|s} \mathbf{1}^T = I
\]
其中 $\alpha_{s|s} = \dfrac{\alpha_s}{\alpha_s} = 1$，$\beta_{s|s} \bar{\pi}_{s|s} = \beta_s \bar{\pi}_s - \dfrac{\alpha_s}{\alpha_s} \beta_s \bar{\pi}_s = 0$

\textbf{② 假设：}假设 $Q_{t|s} = \alpha_{t|s} I + \beta_{t|s} \bar{\pi}_{t|s} \mathbf{1}^T$ 成立。

\textbf{③ 归纳：}证明 $Q_{t+\Delta|s}$ 也满足该形式。
\[
Q_{t+\Delta|s} = \prod_{i=s/\Delta}^{\frac{t+\Delta}{\Delta}-1} \dot{Q}_{\Delta i} = \dot{Q}_t \cdot \prod_{i=s/\Delta}^{\frac{t}{\Delta}-1} \dot{Q}_{\Delta i} = \dot{Q}_t \cdot Q_{t|s}
\]

\[
= (\dot{\alpha}_t I + \dot{\beta}_t \bar{\pi}_t \mathbf{1}^T)(\alpha_{t|s} I + \beta_{t|s} \bar{\pi}_{t|s} \mathbf{1}^T)
\]

由 $\mathbf{1}^T \bar{\pi}_{t|s} \stackrel{H_1}{=} 1$：
\begin{align*}
&= \dot{\alpha}_t \alpha_{t|s} I + \dot{\beta}_t \left\{ \alpha_{t|s} \bar{\pi}_t \mathbf{1}^T + \beta_{t|s} \bar{\pi}_t (\mathbf{1}^T \bar{\pi}_{t|s}) \mathbf{1}^T \right\} + \dot{\alpha}_t \beta_{t|s} \bar{\pi}_{t|s} \mathbf{1}^T \\
&= \dot{\alpha}_t \alpha_{t|s} I + \dot{\beta}_t \left\{ \alpha_{t|s} \bar{\pi}_t \mathbf{1}^T + \beta_{t|s} \bar{\pi}_t \mathbf{1}^T \right\} + \dot{\alpha}_t \beta_{t|s} \bar{\pi}_{t|s} \mathbf{1}^T
\end{align*}

由 $\alpha_{t|s} + \beta_{t|s} \stackrel{H_2}{=} 1$：
\begin{align*}
&= \dot{\alpha}_t \alpha_{t|s} I + \dot{\beta}_t (\alpha_{t|s} + \beta_{t|s}) \bar{\pi}_t \mathbf{1}^T + \dot{\alpha}_t \beta_{t|s} \bar{\pi}_{t|s} \mathbf{1}^T \\
&= \dot{\alpha}_t \alpha_{t|s} I + \dot{\beta}_t \bar{\pi}_t \mathbf{1}^T + \dot{\alpha}_t \beta_{t|s} \bar{\pi}_{t|s} \mathbf{1}^T \\
&= \dot{\alpha}_t \alpha_{t|s} I + (\dot{\beta}_t \bar{\pi}_t + \dot{\alpha}_t \beta_{t|s} \bar{\pi}_{t|s}) \mathbf{1}^T \\
&= \alpha_{t+\Delta|s} I + \beta_{t+\Delta|s} \bar{\pi}_{t+\Delta|s} \mathbf{1}^T
\end{align*}

其中 $\downarrow H_3$：$\alpha_{t+\Delta|s}$，$\downarrow H_4$：$\beta_{t+\Delta|s} \bar{\pi}_{t+\Delta|s}$

\[
Q_{t|s} = \alpha_{t|s} I + \beta_{t|s} \bar{\pi}_{t|s} \mathbf{1}^T \text{ 得证}
\]

注：以上证明仅对离散网格 $s, t \in \{0, \Delta, 2\Delta, \ldots, 1\}$ 成立。下面将其推广到一般的 $s, t \in [0,1]$。

由于 $\alpha_t$，$\bar{\pi}_t$ 可微 $\Rightarrow$
\begin{itemize}[leftmargin=2em]
\item $\alpha_{t+\Delta} = \alpha_t + \Delta \alpha_t' + O(\Delta^2)$，即 $\dfrac{\alpha_{t+\Delta}}{\alpha_t} = 1 + \Delta \dfrac{\alpha_t'}{\alpha_t} + O(\Delta^2)$

于是有：$\lim_{\Delta\to 0} \dfrac{\dot{\alpha}_t - 1}{\Delta} = \dfrac{\alpha_t'}{\alpha_t}$ 存在
\item $\beta_{t+\Delta}\bar{\pi}_{t+\Delta} = \beta_t\bar{\pi}_t + \Delta(\beta_t\bar{\pi}_t)' + O(\Delta^2)$
\item[] 则 $\dot{\beta}_t\dot{\bar{\pi}}_t = \beta_t\bar{\pi}_t + \Delta(\beta_t\bar{\pi}_t)' + O(\Delta^2) - \left(1 + \Delta\dfrac{\alpha_t'}{\alpha_t} + O(\Delta^2)\right)\beta_t\bar{\pi}_t$
\item[] $= \Delta\left\{(\beta_t\bar{\pi}_t)' - \dfrac{\alpha_t'}{\alpha_t}\beta_t\bar{\pi}_t\right\} + O(\Delta^2)$
\item[] 则 $\lim_{\Delta\to 0} \dfrac{\dot{\beta}_t\dot{\bar{\pi}}_t}{\Delta} = (\beta_t\bar{\pi}_t)' - \dfrac{\alpha_t'}{\alpha_t}\beta_t\bar{\pi}_t$ 存在
\end{itemize}

\textbf{则状态转移矩阵存在：}
\[
R_t = \lim_{\Delta\to 0} \frac{Q_{t+\Delta|t} - I}{\Delta} = \lim_{\Delta\to 0} \frac{\dot{Q}_t - I}{\Delta}
= \lim_{\Delta\to 0} \frac{\dot{\alpha}_t - 1}{\Delta} I + \frac{\dot{\beta}_t \dot{\bar{\pi}}_t}{\Delta} \mathbf{1}^T
 = \frac{\alpha_t'}{\alpha_t} I + \left\{(\beta_t\bar{\pi}_t)' - \frac{\alpha_t'}{\alpha_t}\beta_t\bar{\pi}_t\right\}\mathbf{1}^T
\]

\begin{remarkbox}
\refstepcounter{remark}
\textbf{备注 \theremark.} 由 Kolmogorov 前向方程，若状态转移矩阵 $R_t$ 关于 $t$ 连续且 $\sup_{t \in [0,1]} \|R_t\| < \infty$，则存在唯一的转移矩阵族 $\{Q_{t|s}\}_{0 \leq s \leq t \leq 1}$ 满足
\[
\frac{\partial}{\partial t} Q_{t|s} = Q_{t|s} R_t, \quad Q_{s|s} = I
\]
且满足 Chapman-Kolmogorov 方程：对于 $s \leq u \leq t$，
\[
Q_{t|s} = Q_{t|u} Q_{u|s}
\]
\end{remarkbox}

\refstepcounter{proposition}
\begin{propbox}
\textbf{推论 \theproposition（边际转移概率）.}
命题 1.1 中 Markov 链的边际转移概率为
\[
q_t(z_t|x) = \Cat(z_t; Q_t \mathbf{x}), \quad Q_t = \alpha_t I + \beta_t \bar{\pi}_t \mathbf{1}^T
\]
\end{propbox}

\textbf{证明.} 由命题 1.1，取 $s = 0$，则 $Q_t = Q_{t|0}$。由边界条件 $\alpha_0 = 1$，$\beta_0 = 0$，有
\[
\alpha_{t|0} = \frac{\alpha_t}{\alpha_0} = \alpha_t, \quad \beta_{t|0}\bar{\pi}_{t|0} = \beta_t\bar{\pi}_t - \frac{\alpha_t}{\alpha_0}\beta_0\bar{\pi}_0 = \beta_t\bar{\pi}_t
\]
故 $Q_t = \alpha_t I + \beta_t\bar{\pi}_t\mathbf{1}^T$，代入 $q_t(z_t|x) = z_t^T Q_t \mathbf{x}$ 即得
\[
q_t(z_t|x) = \alpha_t z_t^T\mathbf{x} + \beta_t z_t^T\bar{\pi}_t = \alpha_t \delta_{z_t,x} + \beta_t (\bar{\pi}_t)_{z_t}
\]
此即定义 1.3 中的边际前向转移。 \hfill $\square$

\subsection{辅助引理证明}

\refstepcounter{lemma}
\begin{lembox}
\textbf{引理 \thelemma（H1）.} $\mathbf{1}^T \bar{\pi}_{t|s} = 1$
\end{lembox}

\textbf{证明.}
\[
\beta_{t|s} \bar{\pi}_{t|s} = \beta_t\bar{\pi}_t - \frac{\alpha_t}{\alpha_s} \beta_s\bar{\pi}_s
\]
\[
\Rightarrow \beta_{t|s} \mathbf{1}^T\bar{\pi}_{t|s} = \beta_t \mathbf{1}^T\bar{\pi}_t - \frac{\alpha_t}{\alpha_s} \beta_s \mathbf{1}^T\bar{\pi}_s = \beta_t - \frac{\alpha_t}{\alpha_s} \beta_s = \beta_{t|s}
\]
\[
\Rightarrow \mathbf{1}^T\bar{\pi}_{t|s} = 1
\]
\hfill $\square$

\refstepcounter{lemma}
\begin{lembox}
\textbf{引理 \thelemma（H2）.} $\alpha_{t|s} + \beta_{t|s} = 1$
\end{lembox}

\textbf{证明.} $\alpha_{t|s} + \beta_{t|s} = \dfrac{\alpha_t}{\alpha_s} + \beta_t - \dfrac{\alpha_t}{\alpha_s}\beta_s = \dfrac{\alpha_t}{\alpha_s}(1-\beta_s) + \beta_t = \alpha_t + \beta_t = 1$ \hfill $\square$

\refstepcounter{lemma}
\begin{lembox}
\textbf{引理 \thelemma（H3）.} $\alpha_{t+\Delta|s} = \dot{\alpha}_t \alpha_{t|s}$
\end{lembox}

\textbf{证明.}
\[
\alpha_{t+\Delta|s} = \frac{\alpha_{t+\Delta}}{\alpha_s} = \frac{\alpha_{t+\Delta}}{\alpha_t} \cdot \frac{\alpha_t}{\alpha_s} = \dot{\alpha}_t \alpha_{t|s}
\]
\hfill $\square$

\refstepcounter{lemma}
\begin{lembox}
\textbf{引理 \thelemma（H4）.} $\beta_{t+\Delta|s} \bar{\pi}_{t+\Delta|s} = \dot{\beta}_t\bar{\pi}_t + \dot{\alpha}_t\beta_{t|s}\bar{\pi}_{t|s}$
\end{lembox}

\textbf{证明.}
\[
\beta_{t+\Delta|s} \bar{\pi}_{t+\Delta|s} = \beta_{t+\Delta}\bar{\pi}_{t+\Delta} - \frac{\alpha_{t+\Delta}}{\alpha_s} \beta_s\bar{\pi}_s \quad (1)
\]
\[
\dot{\beta}_t\bar{\pi}_t = \beta_{t+\Delta}\bar{\pi}_{t+\Delta} - \frac{\alpha_{t+\Delta}}{\alpha_t} \beta_t\bar{\pi}_t
\]
\[
\dot{\alpha}_t = \frac{\alpha_{t+\Delta}}{\alpha_t}
\]
\[
\beta_{t|s}\bar{\pi}_{t|s} = \beta_t\bar{\pi}_t - \frac{\alpha_t}{\alpha_s}\beta_s\bar{\pi}_s
\]
\[
\Rightarrow \dot{\beta}_t\bar{\pi}_t + \dot{\alpha}_t\beta_{t|s}\bar{\pi}_{t|s} = \beta_{t+\Delta}\bar{\pi}_{t+\Delta} - \frac{\alpha_{t+\Delta}}{\alpha_s}\beta_s\bar{\pi}_s \quad (2)
\]
由 (1) = (2) 得证。 \hfill $\square$

% ========== Section 2 ==========
\begin{center}
{\LARGE\bfseries 2\quad GIDD 的前向转移状态转移矩阵 $R_t$}
\end{center}

\setcounter{section}{2}
\setcounter{subsection}{0}
\setcounter{lemma}{0}

\subsection{前向转移速率}

本节目标：从离散时间的转移矩阵 $Q_{t|s}$ 推导出连续时间的速率矩阵 $R_t$。

\refstepcounter{definition}
\begin{defbox}
\textbf{定义 \thedefinition（CTMC 前向转移）.} 对于起始时刻 $s$ 和终止时刻 $t = s + \Delta$（$\Delta \to 0$），连续时间 Markov 链的前向转移概率满足：
\[
q_{t|s}(z_t|z_s) = \delta_{z_s,z_t} + R_t(z_s, z_t)\Delta + o(\Delta)
\]
其中 $R_t$ 称为\textbf{前向转移速率矩阵}（forward transition rate matrix）。小 $o$ 记号表示渐近次线性项。

速率矩阵 $R_t$ 满足：
\begin{itemize}[leftmargin=2em]
\item 非对角元 $R_t(z, z') \geq 0$（$z \neq z'$）：表示从状态 $z$ 跳转到 $z'$ 的瞬时速率
\item 对角元 $R_t(z, z) = -\sum_{z' \neq z} R_t(z, z')$：保证概率守恒
\end{itemize}
\end{defbox}

\refstepcounter{lemma}
\begin{lembox}
\textbf{引理 \thelemma（Lemma 3.6，GIDD 前向转移速率）.} 设 $t = s + \Delta$，则当 $\Delta \to 0$ 时，命题 1.1 中的转移概率满足：
\[
q_{t|s}(z_t|z_s) = \delta_{z_s,z_t} + R_t(z_s, z_t)\Delta + O(\Delta^2)
\]
其中速率矩阵为
\[
R_t(z_s, z_t) = \frac{\alpha_t'}{\alpha_t}\delta_{z_s,z_t} + z_t^T\left(\beta_t\bar{\pi}_t' - \frac{\alpha_t'}{\alpha_t}\bar{\pi}_t\right)
\]
\end{lembox}

\textbf{证明.} 对 $q_{t|s}(z_t|z_s)$ 作一阶 Taylor 展开。设 $t = s + \Delta$，$\Delta \to 0$。由命题 1.1：
\[
q_{s+\Delta|s}(z_t|z_s) = z_t^T Q_{s+\Delta|s} z_s = z_t^T\left(\alpha_{s+\Delta|s} z_s + \beta_{s+\Delta|s}\bar{\pi}_{s+\Delta|s}\right)
\]

在 $s$ 处展开 $\alpha_{s+\Delta|s}$ 和 $\beta_{s+\Delta|s}\bar{\pi}_{s+\Delta|s}$：
\[
\alpha_{s+\Delta|s} = \frac{\alpha_{s+\Delta}}{\alpha_s} = \frac{\alpha_s + \alpha_s'\Delta + o(\Delta)}{\alpha_s} = 1 + \frac{\alpha_s'}{\alpha_s}\Delta + o(\Delta)
\]
\begin{align*}
\beta_{s+\Delta|s}\bar{\pi}_{s+\Delta|s} &= \beta_{s+\Delta}\bar{\pi}_{s+\Delta} - \frac{\alpha_{s+\Delta}}{\alpha_s}\beta_s\bar{\pi}_s \\
&= \left(\beta_s\bar{\pi}_s + (\beta_s\bar{\pi}_s)'\Delta + o(\Delta)\right) - \left(1 + \frac{\alpha_s'}{\alpha_s}\Delta + o(\Delta)\right)\beta_s\bar{\pi}_s \\
&= \left[(\beta_s\bar{\pi}_s)' - \frac{\alpha_s'}{\alpha_s}\beta_s\bar{\pi}_s\right]\Delta + o(\Delta)
\end{align*}

利用乘积法则简化：
\begin{align*}
(\beta_s\bar{\pi}_s)' - \frac{\alpha_s'}{\alpha_s}\beta_s\bar{\pi}_s &= (1-\alpha_s)'\bar{\pi}_s + \beta_s\bar{\pi}_s' - \frac{\alpha_s'}{\alpha_s}(1-\alpha_s)\bar{\pi}_s \\
&= -\alpha_s'\bar{\pi}_s + \beta_s\bar{\pi}_s' - \frac{\alpha_s'}{\alpha_s}\beta_s\bar{\pi}_s \\
&= \beta_s\bar{\pi}_s' - \alpha_s'\left(1 + \frac{\beta_s}{\alpha_s}\right)\bar{\pi}_s \\
&= \beta_s\bar{\pi}_s' - \frac{\alpha_s'}{\alpha_s}\bar{\pi}_s
\end{align*}

代入 $q_{s+\Delta|s}(z_t|z_s)$：
\begin{align*}
q_{s+\Delta|s}(z_t|z_s) &= z_t^T\left[\left(1 + \frac{\alpha_s'}{\alpha_s}\Delta + o(\Delta)\right)z_s + \left(\beta_s\bar{\pi}_s' - \frac{\alpha_s'}{\alpha_s}\bar{\pi}_s\right)\Delta + o(\Delta)\right] \\
&= \delta_{z_s,z_t} + \underbrace{\left[\frac{\alpha_s'}{\alpha_s}\delta_{z_s,z_t} + z_t^T\left(\beta_s\bar{\pi}_s' - \frac{\alpha_s'}{\alpha_s}\bar{\pi}_s\right)\right]}_{= R_s(z_s,z_t)}\Delta + o(\Delta)
\end{align*}


\textbf{物理意义：}
\begin{itemize}[leftmargin=2em]
\item $R_t(z, z') \, dt$ 表示在 $[t, t+dt)$ 内从状态 $z$ 跳转到 $z'$ 的瞬时概率
\item 对角元 $R_t(z, z) = \dfrac{\alpha_t'}{\alpha_t} < 0$（因 $\alpha_t$ 递减），表示离开状态 $z$ 的速率
\end{itemize}
\hfill $\square$

% ========== Section 3 ==========
\begin{center}
{\LARGE\bfseries 3\quad GIDD 的后向转移速率 $\hat{R}_t$}
\end{center}

\setcounter{section}{3}
\setcounter{subsection}{0}
\setcounter{definition}{0}

本节目标：推导后向转移速率 $\hat{R}_t$，用于从噪声恢复数据。

\subsection{后向转移的定义}

\refstepcounter{definition}
\begin{defbox}
\textbf{定义 \thedefinition（CTMC 后向转移）.} 对于 $s < t$，设 $s = t - \Delta$，$\Delta \to 0$，后向转移概率满足：
\[
p_{s|t}(z_s|z_t) = \delta_{z_t,z_s} + \hat{R}_t(z_t, z_s)\Delta + o(\Delta)
\]
其中 $\hat{R}_t$ 称为后向转移速率。
\end{defbox}

\refstepcounter{lemma}
\begin{lembox}
\textbf{引理 \thelemma（后向转移的贝叶斯形式）.} 真实后向转移概率为：
\[
q_{s|t}(z_s|z_t, x) = \frac{q_{t|s}(z_t|z_s) \cdot q_s(z_s|x)}{q_t(z_t|x)}
\]
\end{lembox}

\textbf{证明.} 由贝叶斯公式：
\[
q_{s|t}(z_s|z_t, x) = \frac{q(z_t, z_s|x)}{q_t(z_t|x)} = \frac{q_{t|s}(z_t|z_s, x) \cdot q_s(z_s|x)}{q_t(z_t|x)}
\]
由马尔可夫性质，给定 $z_s$ 后 $z_t$ 不再依赖于 $x$，即 $q_{t|s}(z_t|z_s, x) = q_{t|s}(z_t|z_s)$，故：
\[
q_{s|t}(z_s|z_t, x) = \frac{q_{t|s}(z_t|z_s) \cdot q_s(z_s|x)}{q_t(z_t|x)}
\]
\hfill $\square$

\textbf{模型后向分布.} 用神经网络 $x_\theta(Z_t, t)$ 预测 $x$ 的分布，定义模型后向分布：
\[
p_\theta(z_s|z_t) := q_{t|s}(z_t|z_s) \cdot \frac{q_s(z_s|x_\theta)}{q_t(z_t|x_\theta)}
\]
其中 $q_t(z_t|x_\theta) := \mathrm{Cat}(z_t; Q_t x_\theta(Z_t, t))$。

\subsection{后向转移速率}

\refstepcounter{lemma}
\begin{lembox}
\textbf{引理 \thelemma（GIDD 后向速率）.} 模型后向分布 $p_\theta(z_s|z_t)$ 的 CTMC 后向速率为：
\[
\hat{R}_t^\theta(z_t, z_s) = -\delta_{z_s,z_t} \sum_{z'} R_t(z', z_t) \frac{q_t(z'|x_\theta)}{q_t(z_t|x_\theta)} + R_t(z_s, z_t)\frac{q_t(z_s|x_\theta)}{q_t(z_t|x_\theta)}
\]
\end{lembox}

\textbf{证明.} 设 $s = t - \Delta$，$\Delta \to 0$。由模型后向分布的定义：
\[
p_\theta(z_s|z_t) = q_{t|s}(z_t|z_s) \cdot \frac{q_s(z_s|x_\theta)}{q_t(z_t|x_\theta)}
\]

代入引理 2.1 的前向转移展开式，并对 $q_s(z_s|x_\theta) = q_{t-\Delta}(z_s|x_\theta)$ 作 Taylor 展开：
\begin{align*}
p_\theta(z_s|z_t) &= \left(\delta_{z_s,z_t} + R_t(z_s, z_t)\Delta + o(\Delta)\right) \cdot \frac{q_t(z_s|x_\theta) - q_t'(z_s|x_\theta)\Delta + o(\Delta)}{q_t(z_t|x_\theta)}
\end{align*}

展开并整理（注意由引理2.1，当$z_s$与$z_t$充分接近时 $\frac{q_t(z_s|x_\theta)}{q_t(z_t|x_\theta)} = 1+O(\Delta)$）：
\begin{align*}
p_\theta(z_s|z_t) &= \delta_{z_s,z_t} + \left[-\delta_{z_s,z_t}\frac{q_t'(z_s|x_\theta)}{q_t(z_t|x_\theta)} + R_t(z_s,z_t)\frac{q_t(z_s|x_\theta)}{q_t(z_t|x_\theta)}\right]\Delta + o(\Delta)
\end{align*}

与后向转移定义 $p_{s|t}(z_s|z_t) = \delta_{z_t,z_s} + \hat{R}_t(z_t, z_s)\Delta + o(\Delta)$ 比较，得：
\[
\hat{R}_t^\theta(z_t, z_s) = -\delta_{z_s,z_t}\frac{q_t'(z_t|x_\theta)}{q_t(z_t|x_\theta)} + R_t(z_s, z_t)\frac{q_t(z_s|x_\theta)}{q_t(z_t|x_\theta)}
\]

利用 Kolmogorov 前向方程 $q_t'(z|x_\theta) = \sum_{z'} q_t(z'|x_\theta) R_t(z', z)$ 消去时间导数：
\[
\hat{R}_t^\theta(z_t, z_s) = -\delta_{z_s,z_t} \sum_{z'} R_t(z', z_t) \frac{q_t(z'|x_\theta)}{q_t(z_t|x_\theta)} + R_t(z_s, z_t)\frac{q_t(z_s|x_\theta)}{q_t(z_t|x_\theta)}
\]
\hfill $\square$












% ========== Section 4 ==========
\begin{center}
{\LARGE\bfseries 4\quad ELBO 下界 (Evidence Lower Bound)}
\end{center}

\setcounter{section}{4}
\setcounter{subsection}{0}
\setcounter{theorem}{0}

\subsection{预备命题}

\refstepcounter{proposition}
\begin{propbox}
\textbf{命题 \theproposition（一般 CT-ELBO，Proposition H.4）.}
对于任意具有边际分布 $q_t(z_t|x)$、前向速率 $R_t(z_s, z_t)$ 和后向速率 $\hat{R}_t(z_t, z_s)$ 的 CTMC 扩散过程，其 CT-ELBO 为：
\[
\log p(x) \geq \mathbb{E}_{t, q_t(z_t|x)} \left[\hat{R}_t(z_t, z_t) - R_t(z_t, z_t) + \sum_{z' \neq z_t} R_t(z', z_t) \frac{q_t(z'|x)}{q_t(z_t|x)} \log \frac{\hat{R}_t(z_t, z') q_t(z_t|x)}{R_t(z', z_t) q_t(z'|x)} \right] + C
\]
其中 $t \sim \mathcal{U}(0,1)$，$C = \mathbb{E}_{q_0(z_0|x)}[\log p(x|z_0)] - \Dkl(q_1(z_1|x) \| p_1(z_1))$ 为 ELBO 常数。
\end{propbox}

\refstepcounter{lemma}
\begin{lembox}
\textbf{引理 \thelemma（权重期望，Lemma H.5）.}
设 $\alpha_t, \beta_t, q_t(z_t|x)$ 如定义 3.1--3.2，权重函数定义为
\[
w_t(z_t, x) = \frac{1}{q_t(z_t|x)} z_t^T \left(\beta_t \bar{\pi}_t' - \frac{\alpha_t'}{\alpha_t}\bar{\pi}_t\right)
\]
则有：
\[
\mathbb{E}_{z_t \sim q_t(\cdot|x)} [w_t(z_t, x)] = -\frac{\alpha_t'}{\alpha_t}
\]
\end{lembox}

\subsection{主要定理}

\refstepcounter{theorem}
\begin{thmbox}
\textbf{定理 \thetheorem（GIDD ELBO，Theorem 3.7）.}
设 $\alpha_t, \beta_t$ 为混合率，$\bar{\pi}_t$ 为混合分布，$q_t(z_t|x)$ 为边际前向分布。定义权重函数：
\[
w_t(z_t, x) = \frac{1}{q_t(z_t|x)} z_t^T \left(\beta_t \bar{\pi}_t' - \frac{\alpha_t'}{\alpha_t}\bar{\pi}_t\right)
\]
则连续时间负 ELBO（CT-NELBO）上界为：
\[
-\log p(x) \leq \mathbb{E}_{t, z_t}\left[w_t(z_t, x) \left(\Dkl(q_t(\cdot|x) \| q_t(\cdot|x_\theta)) + D_{IS}(q_t(z_t|x) \| q_t(z_t|x_\theta))\right)\right] + C
\]
其中 $t \sim \mathcal{U}(0,1)$，$z_t \sim q_t(\cdot|x)$，且 $D_{IS}(p\|q) = \frac{p}{q} - \log\frac{p}{q} - 1$ 为逐点 Itakura-Saito 散度 。
\end{thmbox}

\subsection{定理证明}

\textbf{证明.} 由命题 4.1（Proposition H.4），一般 CT-ELBO 为：
\begin{equation} \label{eq:general_elbo}
\log p(x) \geq \mathbb{E}_{t, q_t(z_t|x)} \left[\hat{R}_t(z_t, z_t) - R_t(z_t, z_t) + \sum_{z' \neq z_t} R_t(z', z_t) \frac{q_t(z'|x)}{q_t(z_t|x)} \log \frac{\hat{R}_t(z_t, z') q_t(z_t|x)}{R_t(z', z_t) q_t(z'|x)}\right] + C
\end{equation}

由引理 3.2（GIDD 后向速率），后向速率为：
\[
\hat{R}_t^\theta(z_t, z_s) = \begin{cases}
R_t(z_s, z_t) \dfrac{q_t(z_s|x_\theta)}{q_t(z_t|x_\theta)} & \text{若 } z_s \neq z_t \\[10pt]
R_t(z_t, z_t) - \displaystyle\sum_{z'} R_t(z', z_t) \dfrac{q_t(z'|x_\theta)}{q_t(z_t|x_\theta)} & \text{若 } z_s = z_t
\end{cases}
\]

将 GIDD 后向速率代入 (\ref{eq:general_elbo})：
\begin{equation} \label{eq:gidd_elbo_intermediate}
\log p(x) \geq \mathbb{E}_{t, z_t} \left[-\sum_{z'} R_t(z', z_t) \frac{q_t(z'|x_\theta)}{q_t(z_t|x_\theta)} + \sum_{z' \neq z_t} R_t(z', z_t) \frac{q_t(z'|x)}{q_t(z_t|x)} \log \frac{q_t(z'|x_\theta) q_t(z_t|x)}{q_t(z_t|x_\theta) q_t(z'|x)}\right] + C
\end{equation}

现在化简期望内的两个求和项。由 GIDD 前向速率（引理 2.1）：
\[
R_t(z_s, z_t) = \frac{\alpha_t'}{\alpha_t}\delta_{z_s,z_t} + z_t^T\left(\beta_t\bar{\pi}_t' - \frac{\alpha_t'}{\alpha_t}\bar{\pi}_t\right)
\]
基于 $w_t(z_t, x)$ 的定义，当 $z' \neq z_t$ 时有 $R_t(z', z_t) = w_t(z_t, x) q_t(z_t|x)$。

\textbf{第一项化简：}
\begin{align*}
\sum_{z'} R_t(z', z_t) \frac{q_t(z'|x_\theta)}{q_t(z_t|x_\theta)} &= \sum_{z'} \left(\frac{\alpha_t'}{\alpha_t}\delta_{z',z_t} + w_t(z_t, x)q_t(z_t|x)\right) \frac{q_t(z'|x_\theta)}{q_t(z_t|x_\theta)} \\
&= \sum_{z'} w_t(z_t, x) q_t(z'|x_\theta) \frac{q_t(z_t|x)}{q_t(z_t|x_\theta)} + \frac{\alpha_t'}{\alpha_t} \frac{q_t(z_t|x_\theta)}{q_t(z_t|x_\theta)} \\
&= w_t(z_t, x) \frac{q_t(z_t|x)}{q_t(z_t|x_\theta)} + \frac{\alpha_t'}{\alpha_t}
\end{align*}

\textbf{第二项化简：}
注意当 $z' \neq z_t$ 时，$\frac{R_t(z', z_t)}{q_t(z_t|x)} = w_t(z_t, x)$，且内部 log 项若 $z' = z_t$ 则为 0。
\begin{align*}
&\sum_{z' \neq z_t} \frac{q_t(z'|x)}{q_t(z_t|x)} R_t(z', z_t) \log \frac{q_t(z'|x_\theta) q_t(z_t|x)}{q_t(z_t|x_\theta) q_t(z'|x)} \\
&= \sum_{z' \neq z_t} w_t(z_t, x) q_t(z'|x) \left( \log \frac{q_t(z'|x_\theta)}{q_t(z'|x)} - \log \frac{q_t(z_t|x_\theta)}{q_t(z_t|x)} \right) \\
&= w_t(z_t, x) \left( \sum_{z'} q_t(z'|x) \log \frac{q_t(z'|x_\theta)}{q_t(z'|x)} + \sum_{z'} q_t(z'|x) \log \frac{q_t(z_t|x)}{q_t(z_t|x_\theta)} \right) \\
&= w_t(z_t, x) \left( -\Dkl(q_t(\cdot|x) \| q_t(\cdot|x_\theta)) + \log \frac{q_t(z_t|x)}{q_t(z_t|x_\theta)} \right)
\end{align*}

\textbf{合并结果：}
将两项代入 (\ref{eq:gidd_elbo_intermediate})：
\begin{align*}
-\log p(x) &\leq -\mathbb{E}_{t, z_t} \left[-w_t(z_t, x) \frac{q_t(z_t|x)}{q_t(z_t|x_\theta)} - \frac{\alpha_t'}{\alpha_t} + w_t(z_t, x) \left( -\Dkl + \log \frac{q_t(z_t|x)}{q_t(z_t|x_\theta)} \right)\right] + C \\
&= \mathbb{E}_{t, z_t} \left[ w_t(z_t, x) \left( \Dkl + \frac{q_t(z_t|x)}{q_t(z_t|x_\theta)} - \log \frac{q_t(z_t|x)}{q_t(z_t|x_\theta)} \right) + \frac{\alpha_t'}{\alpha_t} \right] + C
\end{align*}

由引理 4.1（Lemma H.5）：$\mathbb{E}_{z_t}[w_t(z_t, x)] = -\dfrac{\alpha_t'}{\alpha_t}$，故：
\[
\frac{\alpha_t'}{\alpha_t} = -\mathbb{E}_{z_t}[w_t(z_t, x)]
\]
代入上式，将 $\frac{\alpha_t'}{\alpha_t}$ 吸收进期望内：
\begin{align*}
-\log p(x) &\leq \mathbb{E}_{t, z_t} \left[ w_t(z_t, x) \left( \Dkl + \frac{q_t(z_t|x)}{q_t(z_t|x_\theta)} - \log \frac{q_t(z_t|x)}{q_t(z_t|x_\theta)} - 1 \right) \right] + C \\
&= \mathbb{E}_{t, z_t} \left[ w_t(z_t, x) \left( \Dkl(q_t(\cdot|x) \| q_t(\cdot|x_\theta)) + D_{IS}(q_t(z_t|x) \| q_t(z_t|x_\theta)) \right) \right] + C
\end{align*}
其中 $D_{IS}(p\|q) = \frac{p}{q} - \log\frac{p}{q} - 1$ 为逐点 Itakura-Saito 散度。 \hfill $\square$


% ========== Section 5 ==========
\begin{center}
{\LARGE\bfseries 5\quad 关键引理与详细推导}
\end{center}

\setcounter{section}{5}
\setcounter{subsection}{0}

\subsection{Proposition H.4 的严格证明 (一般 CT-ELBO)}

\textbf{证明.} 扩散 ELBO 的核心量是真实后向转移与模型后向转移之间的 KL 散度。利用 Bayes 公式和前向过程的 Markov 性质 $q_{t|s}(z_t|z_s, x) = q_{t|s}(z_t|z_s)$：
\begin{align}
\Dkl(q_{s|t}(z_s|z_t, x) \| p_{s|t}(z_s|z_t)) &= \sum_{z_s} q_{s|t}(z_s|z_t, x) \log \frac{q_{s|t}(z_s|z_t, x)}{p_{s|t}(z_s|z_t)} \notag \\
&= \sum_{z_s} q_{t|s}(z_t|z_s) \frac{q_s(z_s|x)}{q_t(z_t|x)} \log \frac{q_{t|s}(z_t|z_s) q_s(z_s|x)}{p_{s|t}(z_s|z_t) q_t(z_t|x)} \label{eq:kl-bayes}
\end{align}

为推导连续时间 ELBO，分析当 $\Delta \to 0$（其中 $s = t - \Delta$）时该项的行为。由 CTMC 定义：
\[
\log q_{t|s}(z_t|z_s) = \log(\delta_{z_t, z_s} + \Delta R_s(z_s, z_t) + o(\Delta))
\]

\textbf{Step 1: Taylor 展开}

对于 $z_s = z_t$ 的情形，利用 $\log(1+x) = x - \frac{x^2}{2} + o(x^2)$：
\[
\log q_{t|s}(z_t|z_s) = \log(1 + \Delta R_s(z_t, z_t) + o(\Delta)) = \Delta R_s(z_t, z_t) + o(\Delta)
\]

对于 $z_s \neq z_t$ 的情形，利用 $\Delta \log(\Delta x + o(\Delta)) = \Delta \log \Delta + \Delta \log(x + o(1)) \xrightarrow{\Delta \to 0} \Delta \log x$：
\[
q_{t|s}(z_t|z_s) \log q_{t|s}(z_t|z_s) = [\delta_{z_t,z_s} + \Delta R_s(z_s, z_t) + o(\Delta)] \cdot [\delta_{z_t,z_s} \Delta R_s(z_s,z_t) + (1-\delta_{z_t,z_s}) \log(\Delta R_s(z_s,z_t) + o(\Delta)) + o(\Delta)]
\]

展开后：
\begin{align}
q_{t|s}(z_t|z_s) \log q_{t|s}(z_t|z_s) &= \delta_{z_t,z_s} \Delta R_s(z_s, z_t) + (1-\delta_{z_t,z_s}) \Delta R_s(z_s, z_t) \log R_s(z_s, z_t) + o(\Delta) \label{eq:taylor-q}
\end{align}

类似地：
\begin{align}
q_{t|s}(z_t|z_s) \log p_{s|t}(z_s|z_t) &= \delta_{z_s,z_t} \Delta \hat{R}_s(z_t, z_t) + (1-\delta_{z_s,z_t}) \Delta R_s(z_s, z_t) \log \hat{R}_s(z_t, z_s) + o(\Delta) \label{eq:taylor-p}
\end{align}

以及：
\begin{align}
q_{t|s}(z_t|z_s) \log \frac{q_s(z_s|x)}{q_t(z_t|x)} &= (1-\delta_{z_t,z_s}) \Delta R_s(z_s, z_t) \log \frac{q_s(z_s|x)}{q_t(z_t|x)} + o(\Delta) \label{eq:taylor-ratio}
\end{align}

\textbf{Step 2: KL 散度的 $O(\Delta)$ 展开}

将 \eqref{eq:taylor-q}、\eqref{eq:taylor-p}、\eqref{eq:taylor-ratio} 代入 \eqref{eq:kl-bayes}：
\begin{align}
\Dkl(q_{s|t}(\cdot|z_t,x) \| p_{s|t}(\cdot|z_t)) &= \sum_{z_s} \frac{q_s(z_s|x)}{q_t(z_t|x)} \Big[ q_{t|s}(z_t|z_s) \log q_{t|s}(z_t|z_s) - q_{t|s}(z_t|z_s) \log p_{s|t}(z_s|z_t) \notag \\
&\quad + q_{t|s}(z_t|z_s) \log \frac{q_s(z_s|x)}{q_t(z_t|x)} \Big] \notag \\
&= \underbrace{\Delta R_s(z_t, z_t) - \Delta \hat{R}_s(z_t, z_t)}_{\text{if } z_s = z_t} + \sum_{z_s \neq z_t} \Delta R_s(z_s, z_t) \frac{q_s(z_s|x)}{q_t(z_t|x)} \log \frac{R_s(z_s, z_t) q_s(z_s|x)}{\hat{R}_s(z_t, z_s) q_t(z_t|x)} + o(\Delta) \notag \\
&= \Delta \left( R_s(z_t, z_t) - \hat{R}_s(z_t, z_t) + \sum_{z_s \neq z_t} R_s(z_s, z_t) \frac{q_s(z_s|x)}{q_t(z_t|x)} \log \frac{R_s(z_s, z_t) q_s(z_s|x)}{\hat{R}_s(z_t, z_s) q_t(z_t|x)} \right) + o(\Delta) \label{eq:kl-expansion}
\end{align}

\textbf{Step 3: 从离散到连续}

离散时间扩散 ELBO 为：
\[
\log p(x) \geq -\sum_{i=2}^{T} \mathbb{E}_{q_{\Delta i}(z_t|x)} \left[ \Dkl(q_{\Delta(i-1)|\Delta i}(z_s|z_t, x) \| p_\theta(z_s|z_t)) \right] + C
\]

其中 $C = \mathbb{E}_{q_0(z_0|x)}[\log p(x|z_0)] - \Dkl(q_1(z_1|x) \| p_1(z_1))$。

设 $s = t - \Delta$，取 $\Delta = 1/T \to 0$：
\begin{align*}
\log p(x) &\geq -\sum_{i=2}^{T} \mathbb{E}_{q_{i\Delta}(z_t|x)} \left[ \Dkl(q_{s|t}(z_s|z_t, x) \| p_\theta(z_s|z_t)) \right] + C \\
&= -\underbrace{(1/\Delta - 2)}_{\to 1/\Delta} \mathbb{E}_{t \sim \mathcal{U}\{2\Delta, \ldots, 1\}, q_t(z_t|x)} \left[ \Dkl(\cdots) \right] + C \\
&\xrightarrow{\Delta \to 0} -\frac{1}{\Delta} \mathbb{E}_{t \sim \mathcal{U}(0,1), q_t(z_t|x)} \left[ \Delta \left( R_t(z_t, z_t) - \hat{R}_t(z_t, z_t) + \sum_{z' \neq z_t} R_t(z', z_t) \frac{q_t(z'|x)}{q_t(z_t|x)} \log \frac{R_t(z', z_t) q_t(z'|x)}{\hat{R}_t(z_t, z') q_t(z_t|x)} \right) + o(\Delta) \right] + C
\end{align*}

$\Delta$ 相消，$o(\Delta)/\Delta \to 0$，得到最终的 CT-ELBO：
\[
\log p(x) \geq \mathbb{E}_{t \sim \mathcal{U}(0,1), q_t(z_t|x)} \left[ \hat{R}_t(z_t, z_t) - R_t(z_t, z_t) + \sum_{z' \neq z_t} R_t(z', z_t) \frac{q_t(z'|x)}{q_t(z_t|x)} \log \frac{\hat{R}_t(z_t, z') q_t(z_t|x)}{R_t(z', z_t) q_t(z'|x)} \right] + C
\]
\hfill $\square$

\subsection{Lemma H.5 的严格证明 (权重期望)}

\textbf{证明.} 直接计算 $\mathbb{E}_{z_t}[w_t(z_t, x)]$：
\begin{align}
\mathbb{E}_{z_t}[w_t(z_t, x)] &= \sum_{z_t} q_t(z_t|x) \cdot w_t(z_t, x) \notag \\
&= \sum_{z_t} q_t(z_t|x) \cdot \frac{1}{q_t(z_t|x)} z_t^T \left(\beta_t \bar{\pi}_t' - \frac{\alpha_t'}{\alpha_t}\bar{\pi}_t\right) \notag \\
&= \sum_{z_t} z_t^T \left(\beta_t \bar{\pi}_t' - \frac{\alpha_t'}{\alpha_t}\bar{\pi}_t\right) \label{eq:wt-step1}
\end{align}

由于 $z_t$ 遍历所有 one-hot 向量，$\sum_{z_t} z_t = \mathbf{1}$（全 1 向量），因此：
\begin{align}
\eqref{eq:wt-step1} &= \mathbf{1}^T \left(\beta_t \bar{\pi}_t' - \frac{\alpha_t'}{\alpha_t}\bar{\pi}_t\right) \notag \\
&= \beta_t (\mathbf{1}^T \bar{\pi}_t') - \frac{\alpha_t'}{\alpha_t} (\mathbf{1}^T \bar{\pi}_t) \label{eq:wt-step2}
\end{align}

由于 $\mathbf{1}^T$ 是常向量，可以将其移入时间导数内：
\[
\mathbf{1}^T \bar{\pi}_t' = (\mathbf{1}^T \bar{\pi}_t)'
\]

又因为 $\bar{\pi}_t$ 是概率向量（定义 3.2），对所有 $t$ 有 $\mathbf{1}^T \bar{\pi}_t = 1$，故：
\[
(\mathbf{1}^T \bar{\pi}_t)' = (1)' = 0
\]

代入 \eqref{eq:wt-step2}：
\begin{align}
\mathbb{E}_{z_t}[w_t(z_t, x)] &= \beta_t \cdot 0 - \frac{\alpha_t'}{\alpha_t} \cdot 1 = -\frac{\alpha_t'}{\alpha_t} \notag
\end{align}
\hfill $\square$

\textbf{注.} 这一结果的关键在于 $\bar{\pi}_t$ 始终是概率向量，其分量之和恒为 1，因此对时间的导数为 0。这使得权重函数的期望与具体的混合分布 $\bar{\pi}_t$ 无关，仅依赖于混合率 $\alpha_t$。


% ========== Section 6 ==========
\begin{center}
{\LARGE\bfseries 6\quad 混合调度 (Mixing Schedule) 计算}
\end{center}

\setcounter{section}{6}
\setcounter{subsection}{0}

本章推导特定的混合调度：结合 Masking 和 Uniform noise。
定义 $\alpha_t = \frac{1-t}{C_t}$，$\beta_t \pi_t = \frac{t}{C_t}\mathbf{m} + \frac{c_t}{C_t}\mathbf{u}$，其中 $C_t = 1 + c_t$ 。

\subsection{均匀噪声比例 (Lemma H.7)}

\refstepcounter{lemma}
\begin{lembox}
\textbf{引理 \thelemma（Uniform Noise Ratio, ）.}
设 $c_t = B t^{\gamma/2} (1-t)^{\gamma/2}$。若取 $B = 2^\gamma \frac{p_u}{1-p_u}$，则在 $t=1/2$ 时均匀噪声比例达到最大值 $p_u$。
\end{lembox}

\textbf{证明.}
均匀噪声的比例由 $\frac{c_t}{C_t} = \frac{c_t}{1+c_t}$ 决定。由于 $c_t$ 在 $t=1/2$ 处取得最大值（因为 $t(1-t)$ 对称且在 0.5 最大），故比例也在该处最大。
\[
c_{1/2} = B (1/2)^{\gamma/2} (1/2)^{\gamma/2} = B (1/2)^\gamma = 2^\gamma \frac{p_u}{1-p_u} \cdot 2^{-\gamma} = \frac{p_u}{1-p_u}
\]
则最大比例为：
\[
\frac{c_{1/2}}{1+c_{1/2}} = \frac{\frac{p_u}{1-p_u}}{1 + \frac{p_u}{1-p_u}} = \frac{p_u}{1-p_u + p_u} = p_u
\]
\hfill $\square$

\subsection{权重函数推导 (Lemma H.8)}

\refstepcounter{lemma}
\begin{lembox}
\textbf{引理 \thelemma（GIDD ELBO Weights, ）.}
在此混合调度下，权重函数 $w_t(z_t, x)$ 的闭式解为：
\[
w_t(z_t, x) = z_t^T \left( \frac{\mathbf{m} + (c_t + (1-t)c_t')\mathbf{u}}{(1-t)((1-t)\mathbf{x} + t\mathbf{m} + c_t\mathbf{u})} \right)
\]
\end{lembox}

\textbf{证明.}
首先计算 $c_t'$ ：
\[
c_t' = B \left(\frac{\gamma}{2} t^{\frac{\gamma}{2}-1}(1-t)^{\frac{\gamma}{2}} - \frac{\gamma}{2} t^{\frac{\gamma}{2}}(1-t)^{\frac{\gamma}{2}-1}\right) = \frac{\gamma}{2} \frac{1-2t}{t(1-t)} c_t
\]
计算 $\alpha_t'/\alpha_t$ ：
\[
\frac{\alpha_t'}{\alpha_t} = (\log \alpha_t)' = \left(\log \frac{1-t}{1+c_t}\right)' = -\frac{1}{1-t} - \frac{c_t'}{1+c_t}
\]
计算核心项 $\beta_t \pi_t' - \frac{\alpha_t'}{\alpha_t}\pi_t = \alpha_t (\frac{\beta_t \pi_t}{\alpha_t})'$ ：
\[
\frac{\beta_t \pi_t}{\alpha_t} = \frac{\frac{t\mathbf{m} + c_t\mathbf{u}}{C_t}}{\frac{1-t}{C_t}} = \frac{t\mathbf{m} + c_t\mathbf{u}}{1-t}
\]
求导：
\[
\left(\frac{t\mathbf{m} + c_t\mathbf{u}}{1-t}\right)' = \frac{(1-t)(\mathbf{m}+c_t'\mathbf{u}) - (t\mathbf{m}+c_t\mathbf{u})(-1)}{(1-t)^2} = \frac{\mathbf{m} + (c_t + (1-t)c_t')\mathbf{u}}{(1-t)^2}
\]
最后代入 $w_t(z_t, x)$ 定义：
\[
w_t(z_t, x) = \frac{1}{q_t(z_t|x)} \cdot \alpha_t \cdot \frac{\mathbf{m} + (c_t + (1-t)c_t')\mathbf{u}}{(1-t)^2}
\]
化简即得引理结论。 \hfill $\square$

\end{CJK}
\end{document}
